\scanpage{186}
\mspage{60}
% archive sheet number: 185

dans $J$. (\deleted{Il est immédiat que} \inserted{\unsure{Il s'ensuit|Il s'ensuit donc}} donc que $J$ est \unsure{lui-aussi|lui-même} $\unsure{\pi|\Pi}$-adapté, pour la str. d'ordre induite.) Ceci dit, \deleted{\illegible{1}} l'intersection de deux parties $\pi$-cofinales de $I$ est encore $\pi$-cofinale.

Il suffit de voir qu'elle est cofinale, car sa stabilité par bornes supérieures dans $I$ de parties filtrantes de cardinal est tautologique (si elle est vraie pour des $J_\alpha$, elle est vraie pour $\bigcap_\alpha J_\alpha$). \deleted{Quitte à remplacer} \inserted{Soit} $i_0\in I$, \deleted{$I$ par $I/i_0$, \unsure{on est ramené}} et construisons un élément de $J\cap J'$ qui \deleted{soit $\unsure{\geq|>} i_0$,} majore $i_0$. \deleted{Soit} On construit de proche en proche une \inserted{double} suite

\diagram{\begin{tikzcd}[ampersand replacement=\&,column sep=small,row sep=small]
 i_1\arrow[dr] \& \& i_2\arrow[dr,no head] \& \& \cdots\\
 \& i'_1\arrow[ur] \& \& i'_2\arrow[ur] \&
\end{tikzcd}}

avec les $i_\alpha$ dans $J$, les $i'_\alpha$ dans $J'$, et $i_1\geq i_0$. \deleted{Soit} On a donc

\[
\mathrm{Sup}\,i_\alpha=\mathrm{Sup}\,i'_\alpha\ \deleted{\unsure{\in I}},
\]

et si $i$ est la valeur commune des deux membres, on a $i\in J\cap J'$.

Ceci vu, pour prouver la validité de 5.3, on note que la partie $J'=\mathrm{Ss}^{(M^{\pi}_{\deleted{\unsure{0|\alpha}}})}_{\pi}(X)$ de $I=\mathrm{Ss}_\pi(X)$, formée des $Z_\alpha\subset X$ tels que $Z_\alpha\in M^\pi$, est stable par borne supérieure de parties filtrantes de cardinal $\leq\pi$, vu que $M^\pi$ est stable

